% ============================================================
%  MÓDULO — História, Origem e Ciclos de Evolução
% ============================================================
\chapter{História, Origem e Ciclos de Evolução}

\roteirodidatico
{Uma funcionalidade atual pode ter vindo de uma decisão antiga, ter sido modificada ou estar apenas documentada. História técnica é reconstruir essa sequência, não presumir continuidade.}
{\emph{Commit} registra uma mudança no Git; \emph{spec} descreve intenção e critérios; \emph{ciclo} registra uma etapa de evolução; \emph{proveniência} liga uma afirmação às suas fontes.}
{Comece pelas datas e fontes primárias; relacione cada marco a arquivos e mudanças identificáveis; só então interprete como uma decisão influenciou a arquitetura atual.}
{Uma sequência de versões não comprova causalidade nem sucesso. Para sustentar uma interpretação histórica, indique evidências, lacunas e explicações alternativas.}
{Qual documento ou mudança verificável sustenta cada marco, e que parte da história permanece incerta?}

\casoguiado{reconstituir uma funcionalidade}
{Imagine que uma versão antiga acrescentou uma opção de configuração. A mensagem do commit informa o que mudou; a especificação pode explicar a necessidade; o código mostra a implementação presente; um teste pode indicar o comportamento coberto. Se esses registros não existirem ou não concordarem, descreva a lacuna. A cronologia ajuda a orientar a busca, mas não autoriza inventar a motivação que os registros não documentam.}

\section{De onde veio o Core}

\elemento{Origem: do OpenCode\_Ecosystem ao Core}{\metainfo{ID}{H-01} \hfill
\metainfo{Camada}{0 — Fundação histórica} \hfill
\metainfo{Proveniência}{\pth{CORRIGENDUM.md}}}

\begin{conceito}
Este ecossistema é o \textbf{Core} — a unidade de runtime, memória, qualidade e
orquestração — nascido da unificação e porte do projeto original
\textbf{MarceloClaro/OpenCode\_Ecosystem}. O acervo legado (\pth{specs/})
preserva as especificações numeradas originais (SPEC-001 a SPEC-081+, incluindo o
SPEC-038 do Trust Engine, os SPEC-022 a SPEC-025 da Token Economy e o SPEC-046 da
integração Antigravity), garantindo rastreabilidade histórica das decisões de
design.
\end{conceito}

\begin{descricao}
Os commits mais antigos do repositório atual cuidam do \emph{provisionamento} do
ambiente: detecção de Ubuntu no WSL (evitando \texttt{ERROR\_ALREADY\_EXISTS}),
configuração de \texttt{sudoers} sem senha (instalação não-interativa), dependência
\texttt{zstd} para o Ollama e correção do servidor MCP metacognitive-interconnect.
Ou seja: a história começa com a \emph{infraestrutura funcionando} — e só depois
com o código de negócio.
\end{descricao}

\begin{funcao}
Explicar como decisões anteriores aparecem no repositório atual. Para cada
componente, a relação com uma spec de origem deve ser confirmada nos arquivos; a
numeração histórica, sozinha, não prova que houve porte completo. O
\pth{CORRIGENDUM.md} registra alegações que precisaram ser corrigidas e delimita
o que a evidência disponível permite afirmar.
\end{funcao}

\begin{niveis}
\textbf{Intuição:} um laboratório de ``agentes'' que foi crescendo, cada melhoria
guardando uma explicação escrita do motivo (a ``spec'').\\
\textbf{Implementação:} migração história SPEC$\_{original} \to$ componentes do Core
(\pth{economy/})
com numeração de evolução continuada (R1--R46 $\to$ R47+).\\
\textbf{Formalização:} engenharia de software dirigida a especificação (SPEC-driven) com
cadeia de rastreabilidade origem$\to$implementação$\to$teste, mais autocorreção
institucional de claims (CORRIGENDUM).
\end{niveis}

\section{As duas famílias de especificações}

\begin{resultado}
A pasta \pth{specs/} continha \textbf{371 arquivos Markdown} na conferência de
29/09/2026. As famílias principais incluem:\par
\smallskip

{\footnotesize
\begin{tabularx}{\textwidth}{lYl}
\toprule
Família & Papel & Exemplo \\
\midrule
\textbf{SPEC-001 a SPEC-029} & Especificações nativas do Core (SDD): MetaBus,
Blackboard, Reflexion, Transformer, Orquestrador, Trust Engine, Token Economy,
Scanners, MASWOS, Raciocínio, Ciclos, Integrações CLI, Mapa do Ecossistema &
\pth{mci/metabus.py} \\
\textbf{SPEC-935-R}\cod{*} & Série de evolução contínua do Core (285 arquivos) — regras de
governança, pipeline científico, executores externos, otimizações & \cod{SPEC-935-R432},
\cod{SPEC-935-R435} \\
\textbf{Outras séries SPEC} & Projetos de conteúdo (livros, romances, PDFs, ilustrações)
e integrações específicas & \cod{SPEC-935-R598},
\cod{SPEC-935-R599} \\
\bottomrule
\end{tabularx}
}

\textbf{Fonte:} contagem de arquivos Markdown em \pth{specs/}, 29/09/2026.
\end{resultado}

\begin{aviso}
Nem toda ``spec do acervo legado'' está materializada em código no Core atual —
algumas são portadas em partes (migrations indicam explicitamente o que foi ou não
carregado). Ler a spec não é prova de execução.
\end{aviso}

\clearpage
\section{Ciclos de evolução: o "pulso" do ecossistema}

\elemento{Ciclo de evolução (EvolutionRegistry)}{\metainfo{ID}{H-02} \hfill
\metainfo{Camada}{5 — Evolução} \hfill \metainfo{Rota}{\pth{evolution/cycles.py}}}

\begin{conceito}
O arquivo \pth{evolution/cycles.json} continha \textbf{455 registros} na
conferência de 29/09/2026. Um registro pode descrever uma mudança ou outra
atividade evolutiva; para saber o que ocorreu, é preciso ler seu objetivo, as
alterações e as evidências associadas. O modelo didático deste capítulo organiza
uma mudança em seis etapas: identificar uma necessidade, localizar a solução,
implementá-la, verificar o resultado, registrar a alteração e incorporar o
aprendizado. É um roteiro para entender o processo; não significa que os 455
registros tenham seguido exatamente as mesmas etapas.
\end{conceito}

\begin{metodo}
\textbf{Pipeline de auto-evolução} (inspirado no motor \pth{autoevolve}):
\begin{enumerate}[leftmargin=1.4em,itemsep=1pt]
  \item \textbf{SENSE} — identifica lacuna ou gargalo (ex.: um agente do catálogo
        sem implementação robusta).
  \item \textbf{DISCOVER} — localiza componente/skill catalógico candidato.
  \item \textbf{INSTALL} — integra no \pth{evolution/} / no código do Core.
  \item \textbf{VERIFY} — valida com os testes da família \pth{test_r*}.
  \item \textbf{EVOLVE} — registra o ciclo via \pth{evolution_registry.record(...)}.
  \item \textbf{LEARN} — a memória metacognitiva passa a enxergar a nova
        capacidade.
\end{enumerate}
\end{metodo}

\begin{flowch}[Modelo do ciclo de evolução]{fonte: \pth{evolution/} e motor \pth{autoevolve}; 3 fases por linha}
\matrix[matrix of nodes,name=cy,row sep=4mm,column sep=4mm,nodes={fns},ampersand replacement=\&]{
 |[fne]|SENSE \& |[fns]|DISCOVER \& |[fns]|INSTALL \\
 |[fns]|LEARN \& |[fns]|EVOLVE \& |[fde]|VERIFY (testes) \\
};
\draw[seta] (cy-1-1)--(cy-1-2);
\draw[seta] (cy-1-2)--(cy-1-3);
\draw[seta] (cy-1-3)--(cy-2-3);
\draw[seta] (cy-2-3)--(cy-2-2);
\draw[seta] (cy-2-2)--(cy-2-1);
\draw[seta,plumAccent,dashed] (cy-2-3.south) to[out=-90,in=-90,looseness=1.35]
  node[below=2pt,fill=papel,inner sep=1pt,text=plumAccent,font=\scriptsize]{falhou}
  (cy-1-2.south);
\end{flowch}

\begin{descricao}
\textbf{Como ler o desenho.} As setas ligam \emph{SENSE}, \emph{DISCOVER},
\emph{INSTALL}, \emph{VERIFY}, \emph{EVOLVE} e \emph{LEARN}, nesta ordem.
Assim, uma necessidade identificada leva à busca de uma solução; a solução é
instalada, verificada e registrada; por fim, o aprendizado retorna ao sistema.
A seta tracejada marcada ``falhou'' volta da verificação para a descoberta: um
resultado reprovado exige revisar a solução antes de tentar novamente.
\textbf{Justificativa.} A sequência separa a intenção da mudança, sua execução e
a evidência de que o resultado atende ao critério. Essa separação favorece a
reprodutibilidade porque permite a outra pessoa reconstruir o que foi alterado e
qual evidência sustentou a conclusão. Peng discute a reprodutibilidade como um
padrão mínimo para avaliar alegações científicas quando uma replicação
independente completa não está disponível. A aplicação ao Core é uma analogia
metodológica: o artigo não descreve o registro de ciclos deste repositório
\citref{PENG, R. D. Reproducible research in computational science. \emph{Science}, v.~334, n.~6060, p.~1226-1227, 2011}{10.1126/science.1213847}{A sequência separa intenção, execução e verificação; sua relação com Peng é uma analogia de método. O artigo sustenta a reprodutibilidade como padrão mínimo para avaliar alegações, não uma regra específica de backup ou registro de software.}{Reproducibility has the potential to serve as a minimum standard for judging scientific claims when full independent replication of a study is not possible}{A reprodutibilidade tem o potencial de servir como padrão mínimo para julgar alegações científicas quando a replicação independente completa de um estudo não é possível.}{PENG, 2011, p.~1226--1227}.
\end{descricao}

\begin{funcao}
Transformar ``mudança numérica de versão'' em \emph{memória documentada}: qualquer
agente pode consultar \emph{o que} mudou, \emph{em qual spec} e \emph{com qual
rastreio de teste}.
\end{funcao}

\begin{niveis}
\textbf{Intuição:} um diário de bordo automático — toda melhoria tem página própria.\\
\textbf{Implementação:} registro sequencial R47+, contagem total presente em
\pth{evolution/cycles.json} (455); verificação por família de testes.\\
\textbf{Formalização:} base de dados de traços de evolução (spec $\to$ código $\to$ teste),
permitindo estudo longitudinal de como o ecossistema cresce e autocorrige.
\end{niveis}

\begin{limite}
O registro documenta \emph{o que foi feito}, não \emph{se funcionou fora do
repositório}. Verificação em contexto aberto (benchmark externo) é outra etapa,
sujeita a anti-overclaim.
\end{limite}

\clearpage
\section{Marcos históricos}

\elemento{Linha do tempo resumida}{\metainfo{ID}{H-03} \hfill
\metainfo{Camada}{0 — Timeline} \hfill
\metainfo{Fonte}{git log (371 commits), \pth{CORRIGENDUM.md}}}

\begin{resultado}
\smallskip
{\footnotesize
\begin{tabularx}{\textwidth}{lYl}
\toprule
Fase & Marcos & Evidência \\
\midrule
1. Fundação & Provisionamento WSL+Ollama; correções do MCP metacognitive;
porta Trust Engine, Token Economy, Antigravity & commits iniciais; \pth{SPEC-001..015} \\
2. Metacore & MetaBus, Blackboard A2A, Reflexion, Transformer, Orquestrador,
Protocolo de Agentes & SPEC-001 a SPEC-006 \\
3. Qualidade & Trust Engine, Token Economy, pipeline de Scanners, MASWOS
acadêmico, Raciocínio, Ciclos & SPEC-007 a SPEC-012 \\
4. Governança & SPEC-935-R110 (anti-overclaim); CORRIGENDUM; Scientific
Governance Pipeline (OQS, VSEE, EGS); crítico honesto R142 & \pth{CORRIGENDUM.md},
\pth{mci/metacognitive_evaluator.py} \\
5. Catálogo e skills & 216 agentes; skills especializadas; migração R380
(editores MASWOS); Reversa R471; Landscape R482/R521 & \pth{academic/maswos.py},
\pth{reversa_universal/skill_dispatch.py} \\
6. Executores externos & Goose (R598), Plandex (R599), Gemini CLI (R600),
NotebookLM bridge (R601), Reasonix (R602), Haystack (R604), documentos+doctor
(R603), MiroFish proxy (R605) & \pth{specs/SPEC-935-R59*}
\\
7. Otimização local & R607: perfil mínimo de contexto, \pth{num_ctx=24576},
modelos locais afinados, rage medido & \pth{specs/SPEC-935-R607.md} \\
\bottomrule
\end{tabularx}
}
\end{resultado}

\begin{processo}
\textbf{Como a história é preservada:}
\begin{enumerate}[leftmargin=1.4em,itemsep=1pt]
  \item Specs nativas e série R* em \pth{specs/}.
  \item Registro de ciclos em \pth{evolution/}; consulte o diagnóstico para a contagem atual.
  \item Self-correção em \pth{CORRIGENDUM.md} (ex.: tabela 5 estrelas revisada como
        inventário de recursos; ``Qualis A1'' reescrito como ``meta de rubrica'').
  \item Mapa técnico: \pth{ARCHITECTURE.md}; a leitura guiada deste livro
        acompanha o fluxo nos capítulos de arquitetura e orquestração.
\end{enumerate}
\end{processo}

\begin{aviso}
Datas exatas de cada marcos dependem do histórico de commits remoto; este manual
registra \emph{ordem e provas locais}, não calendário completo de deploy.
\end{aviso}


\section{Resumo e exercícios}

\begin{resultado}
\textbf{O que você deve ter aprendido:} (1) a série \pth{SPEC-935-R*} é um
rastro de especificações ligado à evolução, separado do registro de ciclos;
(2) ciclo de evolução
SENSE/DISCOVER/INSTALL/VERIFY/EVOLVE/LEARN; (3) \pth{migrations/} sinaliza o
que foi ou não portado.
\textbf{Exercício sugerido:} escolha uma spec da série R*, leia o frontmatter e
relacione os critérios de aceitação com uma \pth{BlackboardTask} de prova.
\end{resultado}

% ============================================================

%  Gerado por gen_mod14_ativacao.py (R613) — seção de ativação e cálculo
%  para o módulo mod-01b-historia.tex. Edições manuais são preservadas nas próximas
%  execuções (marcador impede reescrita).
% ============================================================

\section[Leitura guiada]{Leitura guiada — História e evolução: do registro à interpretação}

\elemento{Leitura guiada — História e evolução}{\metainfo{ID}{N-01b} \hfill \metainfo{Ciclo}{R614} \hfill \metainfo{Estilo}{do fato à explicação}}

\begin{descricao}
\textbf{O fenômeno.} Abra os \emph{commits mais antigos} deste repositório. O que
se encontra não é uma ideia brilhante de arquitetura: são tarefas de
\emph{provisionamento} --- detectar Ubuntu no WSL para evitar
\cod{ERROR_ALREADY_EXISTS}, configurar \pth{sudoers} sem senha para instalação não
interativa, instalar \pth{zstd} para o Ollama, consertar o servidor MCP
\pth{metacognitive-interconnect}. Há um fato a explicar: \emph{a história do Core
começa pela infraestrutura, não pelo código de negócio}. Por quê? Porque sem
ambiente reprodutível nenhuma outra camada pode ser medida --- o mesmo motivo pelo
qual um experimento de física só vale depois que o aparato é calibrado.
\end{descricao}

\begin{definicao}
\textbf{Grandezas e definições.}
\begin{itemize}[leftmargin=1.3em,itemsep=1pt]
  \item \textbf{Spec} ($S$): documento que declara \emph{o que} deve existir e \emph{por quê}. Total atual: $S = 371$.
  \item \textbf{Série de evolução} ($R$): sequência numerada $R47, R48, \dots$ que registra mudanças do Core; suas specs somam 285 arquivos da família \cod{SPEC-935-R}.
  \item \textbf{Ciclo} ($C$): registro de uma mudança relevante (SENSE $\to$ LEARN); total $C = 455$.
  \item \textbf{Fração de evolução} ($f$): quanto a série contínua representa do acervo.
\end{itemize}
\end{definicao}

\begin{metodo}
\textbf{Dedução passo a passo.} A rastreabilidade origem $\to$ implementação $\to$
teste não é retórica; é uma cadeia que se pode \emph{contar}.
\begin{enumerate}[leftmargin=1.4em,itemsep=2pt]
  \item Para cada componente, a relação com uma spec de origem precisa ser conferida; a existência de uma spec, sozinha, não demonstra que o componente foi implementado.
  \item Para não perder a origem, a numeração de evolução foi \emph{continuada}: $R1$--$R46$ (legado) $\to$ $R47+$ (Core).
  \item Quando uma mudança não aponta para spec nem teste, parte da trilha de intenção e verificação fica ausente, o que torna a auditoria mais difícil.
  \item Por isso, contagens de specs e de ciclos ajudam a descrever o acervo, mas não substituem a leitura das ligações entre os artefatos.
\end{enumerate}
\end{metodo}

\begin{resultado}
\textbf{Exemplo resolvido (números do repositório).}
\begin{enumerate}[leftmargin=1.4em,itemsep=1pt]
  \item Fração da série contínua: $f = 285/371 \approx 0{,}768$ --- cerca de \textbf{77\%} dos arquivos Markdown de \pth{specs/} pertencem à família \cod{SPEC-935-R}.
  \item Razão entre arquivos de especificação e registros de ciclo: $371/455 \approx 0{,}815$. Essa razão compara dois totais; não demonstra que cada ciclo criou uma spec nem que os conjuntos se correspondam um a um.
\end{enumerate}
\textbf{Leitura correta:} são 371 arquivos Markdown em \pth{specs/}, 285 deles
com o prefixo \cod{SPEC-935-R}, e 455 registros em \pth{evolution/cycles.json}.
Essas contagens são reprodutíveis, mas a auditoria do conteúdo exige ler as
relações entre especificações, mudanças e testes.
\end{resultado}

\begin{resultado}
\textbf{Como interpretar a série temporal.} Em 29/09/2026, o arquivo
\pth{evolution/cycles.json} continha 455 registros com carimbo temporal; o mais
antigo era de 22/07/2026. A diferença entre essas datas é de 69 dias, portanto a
razão aritmética é $455/69 \approx 6{,}6$ registros por dia. Esse valor descreve
somente a distribuição dos registros no arquivo: não mede produtividade,
frequência de mudanças independentes nem qualidade. Registros importados ou
criados em lote alteram a média. Para reproduzir a contagem, leia a chave
\cod{cycles} de \pth{evolution/cycles.json} e conte seus elementos; em seguida,
examine as datas antes de interpretar qualquer taxa.
\end{resultado}

\begin{aviso}
\textbf{Limite de validade.} Contar specs \emph{não} prova que a spec foi
executada: parte do acervo legado foi portada apenas em parte. A verificação
\emph{de execução} é o teste (\pth{tests/test_r*.py}); a verificação \emph{de
mérito} externo continua pendente (anti-overclaim).
\end{aviso}

\begin{resultado}
\textbf{Exercícios.} (1) Se o acervo crescer para 400 specs mantendo a razão,
quantas serão da família \cod{SPEC-935-R}? (2) Explique por que ``a história
começa pela infraestrutura'' é uma afirmação verificável nos commits, e não uma
opinião.
\end{resultado}

% ATIVACAO-R613
\section[Ativação e cálculo]{Ativação e cálculo — História e linhas do tempo}
\elemento{ativ-01b}{\metainfo{Ciclo}{R613} \hfill \metainfo{Módulo}{mod-01b-historia.tex}}

\begin{processo}
GATILHO. consulta cronológica ("quando foi criado", "linha do tempo").
ROTA. navegação temporal direta pelo capítulo.
FUNÇÃO. situar cada componente do ecossistema na linha do tempo da IA e da engenharia de software.
\end{processo}

\begin{resultado}
CÁLCULO. janelas temporais = 4 (pré-1980, 1980--2000, 2000--2020, 2020--hoje); cada janela mapeada a marcos verificáveis.
\end{resultado}

\begin{descricao}
JUSTIFICATIVA TEÓRICA. o campo de agentes apresenta, no próprio roteiro de pesquisa, a organização histórica por conceitos e aplicações. O lastro teórico vem da citação que segue: recorte conferido em 29/09/2026, com a página de origem e tradução livre e fiel.
\end{descricao}

\begin{aviso}
LIMITE E ANTI-OVERCLAIM. linha do tempo é didática; datas de DOI/impressão podem diferir das datas de publicação on-line. A verificação atesta a existência e a
localização da fonte (R110), não o mérito da afirmação.
\end{aviso}

\noindent\textit{Interpretação:} Um roteiro de pesquisa organiza relações entre problemas, mecanismos e aplicações; não funciona como catálogo definitivo de ferramentas nem como prova de que uma implementação particular é superior. Esta distinção orienta a leitura histórica do módulo: a literatura situa as escolhas arquiteturais, enquanto o texto explicita quais delas foram adotadas no Core e quais permanecem decisões locais \citref{JENNINGS, N. R.; SYCARA, K.; WOOLDRIDGE, M. A roadmap of agent research and development. \emph{Autonomous Agents and Multi-Agent Systems}, v.~1, n.~1, p.~7-38, 1998}{10.1023/A:1010090405266}{p.~7--38 (resumo, p.~7). é a moldura histórica e conceitual: o roteiro do campo identifica conceitos-chave e aplicações e como se relacionam — exatamente a função deste módulo para o ecossistema.}{It aims to identify key concepts and applications, and to indicate how they relate to one-another}{O objetivo é identificar conceitos-chave e aplicações e indicar como eles se relacionam entre si.}{JENNINGS, 1998, p.~7}
